\section{Heuristic monomorphisation}
\label{sec:mono}

Monomorphisation replaces (or supplements) premises that quantify
over types by finitely many \emph{instances} at ground types.  We
formulate it as an operation on $\CICm$ propositions --- instantiate,
normalise, retranslate --- rather than as an operation on the
generated first-order formulas.  The choice matters: after
instantiation, $\beta\iota\delta$-redexes disappear
(\cref{ex:mono-beta}), the type-directed translation re-runs on the
instance and produces structurally better formulas, and the
soundness argument is a two-line application of the semantics.  A
purely first-order instantiation of the translated generic formula
would have to reproduce these effects through bridging axioms, which
the design constraint of a \emph{shallow} translation --- every
transformation is a per-occurrence formula manipulation, with no
mediating axiomatic theory --- rules out.

Throughout this section, fix an admissible environment $\Env$ and a
standard model $\MM$.

\subsection{Instances of premises}

\begin{definition}[Prefix; instantiation]
\label{def:instantiation}
Let $\varphi$ be a closed proposition in $\beta\iota\delta$-normal
form, written with maximal product prefix
\[
  \varphi \;=\; \Pi x_1{:}B_1.\ \cdots\ \Pi x_n{:}B_n.\ \psi
  \qquad (\psi \text{ not a product}).
\]
An \emph{instantiation} of $\varphi$ is a partial function $\theta$
on $\{x_1, \ldots, x_n\}$ such that, processing positions in order,
each $x_i \in \dom\theta$ is mapped to a \emph{closed}
$\beta\iota\delta$-normal term $\theta(x_i)$ with $\vdash \theta(x_i)
: B_i\theta^{<i}$, where $\theta^{<i}$ denotes the simultaneous
substitution of the already-processed values.  The \emph{instance}
is
\[
  \inst{\varphi}{\theta} \;\coloneqq\;
  \mathrm{nf}\Bigl( \Pi_{\,i \notin \dom\theta}\, x_i :
  B_i\theta.\ \psi\theta \Bigr),
\]
the $\beta\iota\delta$-normal form of the result of deleting the
instantiated binders and substituting throughout (no capture can
occur: the substituted terms are closed).  We call $\theta$ a
\emph{type instantiation} if every $\theta(x_i)$ is a data type
(then necessarily $B_i\theta^{<i} = \STypeI$); the general case
additionally instantiates data and kind-level binders by closed
terms and is used for higher-order instance generation
(\cref{rem:tier2}).  A closed proposition is \emph{monomorphic} if
$\STypeI$ does not occur in it.
\end{definition}

\begin{lemma}[Instances are well-formed]
\label{lem:inst-typed}
$\inst{\varphi}{\theta}$ is a closed proposition, i.e.\
$\vdash \inst{\varphi}{\theta} : \SProp$.
\end{lemma}

\begin{proof}
By induction on $n$.  If $x_1 \notin \dom\theta$, apply the
induction hypothesis to $\psi' = \Pi x_2 \ldots \psi$ in the context
$x_1{:}B_1$ (the definition and the claim relativise to a context of
uninstantiated binders in the obvious way) and rebuild the product
with the rule that formed it.  If $x_1 \in \dom\theta$ with value
$\tau$: by generation, $x_1{:}B_1 \vdash \Pi x_2 \ldots \psi :
\SProp$; by the substitution property \cref{prop:metatheory}(2)
(in the form of \cref{lem:sort-invariance} when $B_1 = \STypeI$;
in general by \cref{prop:metatheory}(2) directly), $\vdash (\Pi x_2
\ldots \psi)[\tau/x_1] : \SProp$, and we proceed by the induction
hypothesis.  Normalisation preserves typing by subject reduction
\cref{prop:metatheory}(3), and the normal form exists and is unique
by \cref{prop:metatheory}(1),(5).
\end{proof}

\begin{lemma}[Semantic instantiation]
\label{lem:inst-valid}
If $\varphi$ is valid in $\MM$ then so is $\inst{\varphi}{\theta}$.
\end{lemma}

\begin{proof}
By induction on $n$; we show the stronger contextual statement: if
$\Gamma$-proposition $\Pi x{:}B.\,\chi$ has
$\semr{\Pi x{:}B.\,\chi}{\rho} = \holds$ and $\vdash \tau : B$ is
closed, then $\semr{\chi[\tau/x]}{\rho} = \holds$; and if $x$ is not
instantiated, validity passes to all $\semr{\chi}{\rho[x \mapsto
a]}$, $a \in \semr{B}{\rho}$, and is recovered for the rebuilt
product.  Both parts are immediate from \ref{M:pi}: validity of the
product means $\semr{\chi}{\rho[x \mapsto a]} = \holds$ for
\emph{every} $a \in \semr{B}{\rho}$; for the instantiated case take
$a = \sem{\tau}$, which lies in $\semr{B}{\rho}$ by \ref{M:typing},
and use \ref{M:subst} to identify
$\semr{\chi[\tau/x]}{\rho} = \semr{\chi}{\rho[x \mapsto
\sem{\tau}]}$.  Finally, normalisation does not change denotations
by \ref{M:conv}.
\end{proof}

Instances are also \emph{derivable} from their premises inside the
type theory, which matters for the reconstruction phase: any proof
using an instance can be replayed from the original premise.

\begin{proposition}[Syntactic instantiation]
\label{prop:inst-derivable}
If $\vdash p : \varphi$ then there is a term $p^{\theta}$ with
$\vdash p^{\theta} : \inst{\varphi}{\theta}$.
\end{proposition}

\begin{proof}
By induction on the prefix, define $p^{\theta}$: for $x_1 \in
\dom\theta$ take $(p\,\theta(x_1))^{\theta'}$ (application rule,
then the induction hypothesis for the substituted suffix); for $x_1
\notin \dom\theta$ take $\lambda x_1{:}B_1\theta.\ (p\,x_1)^{\theta}$
(abstraction rule).  The resulting type is convertible to
$\inst{\varphi}{\theta}$, and the conversion rule concludes.
\end{proof}

\subsection{Monomorphisation of problems}

\begin{definition}[Monomorphisation]
\label{def:mono-op}
Let $(P, \varphi_G)$ be a problem (\cref{def:transl-problem}).
Recall that the emitted inductive axioms (D3) are the
$\mathsf{F}$-translations of closed $\CICm$ propositions; we treat
these propositions as (implicit, valid) premises, so that
instantiation applies to them uniformly.  A \emph{monomorphisation}
is a family $\Theta = (\Theta_p)_{p}$ assigning to each premise $p$
(explicit or implicit) a finite set of instantiations of its
statement $\varphi_p$.  It induces two premise sets:
\begin{align*}
  \Mono^{\mathrm{keep}}_{\Theta}(P) &\;=\;
    P \;\cup\; \setof{\inst{\varphi_p}{\theta}}{p, \ \theta \in \Theta_p},\\
  \Mono^{\mathrm{drop}}_{\Theta}(P) &\;=\;
    \setof{\inst{\varphi_p}{\theta}}{p, \ \theta \in \Theta_p,\
      \inst{\varphi_p}{\theta} \text{ monomorphic}}
    \;\cup\; \setof{\varphi_p}{p, \ \varphi_p \text{ monomorphic}}.
\end{align*}
The monomorphised problems are $\mathrm{Transl}(\Mono_\Theta(P),
\varphi_G)$ for either policy.
\end{definition}

\begin{theorem}[Soundness of monomorphisation]
\label{thm:mono-sound}
For every monomorphisation $\Theta$ and either policy, the
monomorphised problem is $\MM$-sound (\cref{def:si}).  Consequently,
if the ATP proves the translated goal from the monomorphised
premises, the goal is valid in $\MM$.
\end{theorem}

\begin{proof}
Every element of $\Mono_\Theta(P)$ is either a premise statement or
an instance of one, hence valid in $\MM$ by hypothesis and
\cref{lem:inst-valid}; the implicit (D3) premises are valid by
\ref{M:ind} as in \cref{thm:truth}.  \Cref{thm:base-soundness}
applied to the premise set $\Mono_\Theta(P)$ gives the claim.  (The
axiom store grows --- translating instances lifts new subterms ---
but \cref{thm:truth} covers all emitted axioms.)
\end{proof}

Note what this theorem does \emph{not} depend on: how $\Theta$ was
found.  Any instance-selection heuristic whatsoever --- including an
unsound-looking syntactic matcher --- yields a sound pipeline,
because soundness is carried entirely by
\cref{lem:inst-valid,lem:inst-typed}.  Heuristics affect only
effectiveness.  This modularity is the main payoff of formulating
monomorphisation at the $\CICm$ level.

\begin{proposition}[Conservativity of \emph{keep}]
\label{prop:mono-conservative}
Every first-order proof available before monomorphisation remains
available after $\Mono^{\mathrm{keep}}_{\Theta}$: if
$\mathrm{Transl}(P, \varphi_G) = (\Phi, \gamma)$ proves $\gamma$,
then the keep-policy problem $(\Phi', \gamma)$ also proves
$\gamma$.
\end{proposition}

\begin{proof}
Translation is monotone in the premise set and idempotent on the
axiom store (lifted constants and their axioms are indexed by
$\alpha$-classes, \cref{def:lifted}), so $\Phi \subseteq \Phi'$, and
classical provability is monotone.
\end{proof}

\begin{proposition}[Incompleteness of \emph{drop}]
\label{prop:drop-incomplete}
There is a problem $(P, \varphi_G)$ and a monomorphisation $\Theta$
such that $\mathrm{Transl}(P, \varphi_G)$ proves its conjecture but
the drop-policy problem does not.
\end{proposition}

\begin{proof}
Let $\Env$ declare data types $\sigma, \tau$ with closed inhabitants
$a : \sigma$, and $f : \Pi A{:}\STypeI.\,A \to A$, and take the
premise $\varphi_p = \Pi A{:}\STypeI.\,\Pi x{:}A.\ f\,A\,x \ceq_A
x$ with goal $\varphi_G = f\,\sigma\,a \ceq_\sigma a$, and
$\Theta_p = \{[A \mapsto \tau]\}$.  The generic problem proves the
goal (instantiate the guarded universal with $\hat\sigma$ and
$\hat{a}$; the guards $\hastype(\hat\sigma, \lift{\STypeI})$ and
$\hastype(\hat{a}, \hat\sigma)$ are typing atoms of the
translation).  For the drop-policy problem, only the
$\tau$-instance is present; a countermodel is obtained from
$\FMstar$ by reinterpreting $\ap$ at the point
$(\sem{f}, \sem{\sigma})$: replace $\sem{f}$'s value at $\sigma$
by a non-identity function $g$ on $\sem{\sigma}$ with $g(\sem{a})
\neq \sem{a}$ (possible whenever $|\sem{\sigma}| \geq 2$; formally,
reinterpret the constant $\hat{f}$ by a graph agreeing with
$\sem{f}$ except at argument $\sem{\sigma}$).  All axioms of the
drop-policy problem mention $f$ only at $\tau$ or generically
through typing atoms that remain true (the modified graph still
maps $\sem{\sigma}$ into $\sem{\sigma} \to \sem{\sigma}$), while
the conjecture is now false.
\end{proof}

\begin{example}[Why instances are built before translation]
\label{ex:mono-beta}
Consider the polymorphic inversion premise for the inductive
predicate $\mathsf{Forall} : \Pi A{:}\STypeI.\,(A \to \SProp) \to
\mathsf{list}\,A \to \SProp$:
\[
  \varphi \;=\; \forall A{:}\STypeI.\ \forall P{:}A \to \SProp.\
  \forall l{:}\mathsf{list}\,A.\
  \mathsf{Forall}\,A\,P\,l \leftrightarrow
  \bigl( l \ceq \mathsf{nil}\,A \;\vee\; \cexists{x}{A}
  \cexists{l'}{\mathsf{list}\,A}\, (P\,x \wedge \cdots) \bigr).
\]
Its translation contains the higher-order application pattern
$\PR(P \app x)$ with $P$ a quantified variable, and guards
$\hastype(P, \trCn{A \to \SProp})$ through a lifted type constant
--- both notoriously hard for ATPs.  Suppose the goal concerns
$\mathsf{Forall}\,\mathsf{nat}\,(\lambda x{:}\mathsf{nat}.\,
\mathsf{lt}\,\mathsf{O}\,x)\,l$.  The occurrence yields the
instantiation $\theta = [A \mapsto \mathsf{nat},\ P \mapsto \lambda
x.\,\mathsf{lt}\,\mathsf{O}\,x]$ (a general instantiation in the
sense of \cref{def:instantiation}); in
$\inst{\varphi}{\theta}$ the redex $(\lambda
x.\,\mathsf{lt}\,\mathsf{O}\,x)\,x$ is \emph{normalised away}, so
the instance translates to a first-order formula whose corresponding
subformula is the direct atom
$\PR(\widehat{\mathsf{lt}} \app \hat{\mathsf{O}} \app x)$ ---
first-order, guard-light, and after
\cref{sec:specialization} simply $\mathsf{lt}(\mathsf{O}, x)$.  No
first-order manipulation of the translated generic formula can
achieve this: the $\beta$-step happens in the type theory.
\end{example}

\begin{remark}[Higher-order instances]
\label{rem:tier2}
\Cref{lem:inst-valid,lem:inst-typed,prop:inst-derivable} were
deliberately proved for \emph{general} instantiations, not just type
instantiations.  Instantiating a higher-order premise variable
$P : A \to \SProp$ with a closed abstraction found at an occurrence
(as in \cref{ex:mono-beta}) is therefore covered by the same
soundness theorem.  The practical iteration scheme below applies
unchanged; only the occurrence-matching heuristic differs.  This is
the analogue of Lean-auto's instantiation of ``constant instances
with dependent arguments'' \cite{QianEtAl2025}.
\end{remark}

\subsection{Instance discovery and bounds}
\label{sec:saturation}

Soundness being heuristic-independent, we specify the discovery
algorithm only to the degree needed to state termination and cost
bounds.  The scheme is the standard goal-seeded saturation of
Sledgehammer \cite[\S7.1]{Blanchette2016}, transplanted to constant
instances in the sense of \cref{def:instantiation}.

\begin{definition}[Occurrence extraction; saturation]
\label{def:saturation}
For a normal term $t$, let $\TArgs(t)$ be the finite set of pairs
$(c, \tup{\tau})$ such that an application $c\,\tau_1 \cdots \tau_j$
occurs in $t$ where $c$ is a constant whose type has a product
prefix beginning with $j$ binders of domain $\STypeI$ and each
$\tau_i$ is closed and normal.  Given a problem $(P, \varphi_G)$ and
parameters $K, \Delta_{\mathrm{tot}}, \Delta_c \in \mathbb{N}$, the
saturation computes, for $r = 0, 1, \ldots, K$: the \emph{active
instance pool} $W_r$ ($W_0 = \TArgs(\varphi_G)$), and for each
generic premise $p$ the instantiations obtained by matching the
type-variable prefix of $\varphi_p$ against elements of $W_r$
occurring in $\varphi_p$'s own constant applications; the resulting
instances contribute their $\TArgs$ to $W_{r+1}$.  Generation stops
after round $K$, or as soon as the number of instances generated for
one premise exceeds $\Delta_c$ or the total exceeds
$\Delta_{\mathrm{tot}}$ (excess instances are discarded by a fixed
fair strategy).
\end{definition}

\begin{proposition}
\label{prop:saturation-bounds}
The saturation terminates, produces at most
$\Delta_{\mathrm{tot}}$ instances, and every produced $\theta$ is an
instantiation in the sense of \cref{def:instantiation}; hence
\cref{thm:mono-sound} applies.
\end{proposition}

\begin{proof}
Termination and the cardinality bound are immediate from the caps.
Well-formedness of the produced $\theta$: the matched type arguments
are closed normal subterms of well-typed terms at $\STypeI$-typed
positions, so each binding satisfies the typing requirement of
\cref{def:instantiation} (formally: a subterm occurrence of a
well-typed term is well-typed in its local context; a closed subterm
at a position of type $\STypeI$ is a closed data type).
\end{proof}

Sledgehammer operates with $K = 3$ rounds and a total cap of $200$
new formulas by default, values twice validated experimentally
\cite{Blanchette2016,Desharnais2022}; Lean-auto uses a comparable
global cap \cite{QianEtAl2025}.  We adopt them as defaults.

\subsection{No complete instantiation strategy exists}
\label{sec:undecidability}

Heuristic monomorphisation with the \emph{drop} policy is
incomplete (\cref{prop:drop-incomplete}); with the \emph{keep}
policy it is conservative but the generic formulas remain a burden.
One might hope for a clever computable strategy producing, for each
problem, a finite instance set that is \emph{complete} --- one that
loses nothing.  For polymorphic first-order logic,
Bobot and Paskevich proved this hope vain
\cite[Thm.~1]{BobotPaskevich2011}.  We now prove the analogous
result for $\CICm$, adapting their reduction (from ground
reducibility in combinatory logic, with combinator terms reflected
at the type level) to the dependent-type-theoretic setting.  The
notion of consequence is semantic, matching our soundness
framework.

\begin{definition}[Consequence]
\label{def:consequence}
Let $\Env$ be an environment and $P$ a finite set of closed
propositions, $\varphi$ a closed proposition.  Write $P
\folmodels_{\Env} \varphi$ if for every standard model $\MM$ of
$\Env$ in which all members of $P$ are valid, $\varphi$ is valid in
$\MM$.  (Standard models differ in their interpretation of
parameter constants; \ref{M:ind} pins inductive types only up to
the listed properties.)
\end{definition}

By \cref{lem:inst-valid}, instances are consequences:
$\setof{\inst{\varphi_p}{\theta}}{\ldots} \folmodels_\Env \varphi$
implies $P \folmodels_\Env \varphi$.  A \emph{complete instantiation
strategy} would be a computable converse:

\begin{theorem}[Undecidability of complete monomorphisation]
\label{thm:undecidable}
There is no algorithm that, given an environment $\Env$, a finite
premise set $P$ and a closed goal $\varphi$, computes a finite
monomorphisation $\Theta$ such that
\[
  P \folmodels_{\Env} \varphi
  \quad\Longrightarrow\quad
  \setof{\inst{\varphi_p}{\theta}}{p,\ \theta \in \Theta_p}
  \folmodels_{\Env} \varphi .
\]
\end{theorem}

\begin{proof}
We construct a family of problems encoding weak reduction in
combinatory logic (CL).  CL terms are $w \Coloneqq K \mid S \mid
w\,w$; weak reduction is generated by $K\,x\,y \to x$,
$S\,x\,y\,z \to x\,z\,(y\,z)$ and congruence, and $\to^{*}$ is its
reflexive--transitive closure.  The relation $w \to^{*} w'$ between
ground CL terms is undecidable: by combinatory completeness
\cite{Barendregt1984} there is, for every partial recursive $f$, a
CL term $F$ with $F\,\overline{n} \to^{*} \overline{f(n)}$ for all
$n$ in the domain of $f$ (numerals $\overline{n}$ in normal form),
and $F\,\overline{n} \to^{*} \overline{0}$ is then equivalent to
$f(n) = 0$, which for suitable $f$ is $\Sigma_1$-complete.

\emph{The environment $\Env_{\mathrm{CL}}$.}  Declare data
inductives: $\mathsf{K}_0$ with one $0$-ary constructor
$\mathsf{k}_0$; $\mathsf{S}_0$ with \emph{two} $0$-ary constructors
$\mathsf{s}_0, \mathsf{s}_1$ (the second is a tag ensuring
$\sem{\mathsf{K}_0} \neq \sem{\mathsf{S}_0}$ in every model, by
\ref{M:ind}(b--d): one denotation is a singleton, the other a
two-element set); and $\mathsf{A}$ with parameters $\alpha,
\beta{:}\STypeI$ and one constructor $\mathsf{a}_0 : \Pi
\alpha\beta{:}\STypeI.\ \alpha \to \beta \to
\mathsf{A}\,\alpha\,\beta$.  Declare a parameter
$R : \Pi \alpha\beta{:}\STypeI.\ \alpha \to \beta \to \SProp$.
Reflect CL terms as closed types and canonical points:
\begin{gather*}
  \lift{K} = \mathsf{K}_0,\quad
  \lift{S} = \mathsf{S}_0,\quad
  \lift{w\,w'} = \mathsf{A}\,\lift{w}\,\lift{w'};\\
  p_K = \mathsf{k}_0,\quad
  p_S = \mathsf{s}_0,\quad
  p_{w\,w'} = \mathsf{a}_0\,\lift{w}\,\lift{w'}\,p_w\,p_{w'} .
\end{gather*}
The premise set $P_{\mathrm{CL}}$ consists of the following closed
propositions (all prefixes over $\STypeI$; element variables typed
by the type variables):
\begin{align}
& \forall \alpha.\ \forall u{:}\alpha.\ R\,\alpha\,\alpha\,u\,u
  \tag{refl}\\
& \forall \alpha\beta\gamma.\ \forall u{:}\alpha\, v{:}\beta\,
  w{:}\gamma.\ R\,\alpha\,\beta\,u\,v \to R\,\beta\,\gamma\,v\,w \to
  R\,\alpha\,\gamma\,u\,w
  \tag{trans}\\
& \forall \alpha\beta\gamma.\ \forall u{:}\alpha\, v{:}\beta\,
  w{:}\gamma.\ R\,\alpha\,\beta\,u\,v \to
  R\,(\mathsf{A}\alpha\gamma)\,(\mathsf{A}\beta\gamma)\,
    (\mathsf{a}_0\alpha\gamma\,u\,w)\,(\mathsf{a}_0\beta\gamma\,v\,w)
  \tag{congL}\\
& \forall \alpha\beta\gamma.\ \forall u{:}\alpha\, v{:}\beta\,
  w{:}\gamma.\ R\,\alpha\,\beta\,u\,v \to
  R\,(\mathsf{A}\gamma\alpha)\,(\mathsf{A}\gamma\beta)\,
    (\mathsf{a}_0\gamma\alpha\,w\,u)\,(\mathsf{a}_0\gamma\beta\,w\,v)
  \tag{congR}\\
& \forall \alpha\beta.\ \forall u{:}\alpha\, v{:}\beta.\
  R\,\bigl(\mathsf{A}(\mathsf{A}\mathsf{K}_0\alpha)\beta\bigr)\,
    \alpha\,
    \bigl(\mathsf{a}_0\,(\mathsf{A}\mathsf{K}_0\alpha)\,\beta\,
      (\mathsf{a}_0\mathsf{K}_0\alpha\,\mathsf{k}_0\,u)\,v\bigr)\,
    u
  \tag{redK}\\
& \forall \alpha\beta\gamma.\ \forall u{:}\alpha\, v{:}\beta\,
  w{:}\gamma. \notag\\
& \qquad
  R\,\bigl(\mathsf{A}(\mathsf{A}(\mathsf{A}\mathsf{S}_0\alpha)\beta)\gamma\bigr)\,
    \bigl(\mathsf{A}(\mathsf{A}\alpha\gamma)(\mathsf{A}\beta\gamma)\bigr)\,
    t_S\; t'_S
  \tag{redS}
\end{align}
where $t_S$ and $t'_S \coloneqq
\mathsf{a}_0\,(\mathsf{A}\alpha\gamma)(\mathsf{A}\beta\gamma)\,
(\mathsf{a}_0\alpha\gamma\,u\,w)\,(\mathsf{a}_0\beta\gamma\,v\,w)$
are the evident points of the redex and reduct types (both built
from $\mathsf{a}_0, \mathsf{s}_0, u, v, w$).  The goal for ground CL
terms $w, w'$ is $g(w, w') \coloneqq
R\,\lift{w}\,\lift{w'}\,p_w\,p_{w'}$.

\emph{Claim 1: $P_{\mathrm{CL}} \folmodels_{\Env} g(w,w')$ iff
$w \to^{*} w'$.}

($\Leftarrow$)  Let $\MM$ be any standard model validating
$P_{\mathrm{CL}}$.  By induction on the derivation of $w \to^{*}
w'$: each single step and each closure rule corresponds to one
premise; by \cref{lem:inst-valid} the premise may be semantically
instantiated at the reflected types and canonical points of the
terms involved, and validity of an implication together with
validity of its antecedents yields validity of its consequent
(\ref{M:pi}, \ref{M:conn}).  The base facts are (redK), (redS),
(refl); the guards are satisfied because $\sem{p_v} \in
\sem{\lift{v}}$ by \ref{M:typing}.

($\Rightarrow$)  Suppose $w \not\to^{*} w'$; we exhibit a standard
model validating $P_{\mathrm{CL}}$ but not $g(w,w')$.  Take the
\emph{canonical tagged model}: constructor values are tuples tagged
by the constructor name, e.g.\ $\sem{\mathsf{a}_0}(a, b)(x, y) =
(\texttt{a0}, a, b, x, y)$ (with $\texttt{a0}$ a fixed set distinct
from the other tags); then \ref{M:ind} holds, and for every set $d$
we can define a syntax tree $\mathrm{tr}(d)$ over the alphabet
$\{K, S\} \cup \{\text{atoms}\}$ by $\in$-recursion:
$\mathrm{tr}(d) = K$ if $d = \sem{\mathsf{K}_0}$; $\mathrm{tr}(d) =
S$ if $d = \sem{\mathsf{S}_0}$; $\mathrm{tr}(d) =
\mathrm{tr}(a)\cdot\mathrm{tr}(b)$ if $d = \sem{\mathsf{A}}(a,b)$
for (necessarily unique, by the tagging) $a, b$; and
$\mathrm{tr}(d) = d$ (an atom) otherwise.  Uniqueness of the
decomposition holds because $\sem{\mathsf{A}}(a,b)$ is the set of
tagged tuples $(\texttt{a0},a,b,x,y)$ with $x \in a$, $y \in b$,
from which $a$ and $b$ are recoverable when the set is nonempty, and
which is empty iff $a$ or $b$ is empty --- in the degenerate empty
cases let $\mathrm{tr}$ take the atom branch; reflected types
$\lift{v}$ are nonempty (they contain $\sem{p_v}$) so
$\mathrm{tr}(\sem{\lift{v}}) = v$ for all ground CL terms $v$, by
induction (for the base cases: $\sem{\mathsf{K}_0}$ and
$\sem{\mathsf{S}_0}$ are distinguishable by cardinality).  Extend
weak reduction to trees with atoms (atoms are inert constants).
Interpret
\[
  \sem{R}(a, b)(x, y) \;=\;
  \begin{cases}
    \holds & \text{if } \mathrm{tr}(a) \to^{*} \mathrm{tr}(b),\\
    \emptyset & \text{otherwise}
  \end{cases}
\]
(element arguments ignored).  All six premises are valid: (refl) and
(trans) because $\to^{*}$ is a preorder on trees; (congL), (congR)
because $\mathrm{tr}(\sem{\mathsf{A}}(a,b)) =
\mathrm{tr}(a)\cdot\mathrm{tr}(b)$ whenever the guard-supplied
elements witness nonemptiness of $a$ and $b$ (they do: the premise
hypotheses provide $u \in a$ etc.); (redK), (redS) because
$K\,x\,y \to x$ and $S\,x\,y\,z \to x z (y z)$ hold for arbitrary
trees $x, y, z$.  The goal is not valid:
$\mathrm{tr}(\sem{\lift{w}}) = w \not\to^{*} w' =
\mathrm{tr}(\sem{\lift{w'}})$.

\emph{Claim 2: for every finite monomorphisation $\Theta$ of
$P_{\mathrm{CL}}$, the relation ``\,$\mathrm{Inst}_\Theta
\folmodels_{\Env} g(w, w')$\,'' is decidable, uniformly in
$(\Theta, w, w')$.}  Let $T_0$ be the finite set of closed types
occurring in the instances of $\Theta$ together with $\lift{w},
\lift{w'}$, closed under syntactic subterms.  Emptiness of the
denotation of a closed data type is uniform across standard models
and decidable (a type $\mathsf{A}\,\sigma\,\rho$ is inhabited iff
both arguments are; $\mathsf{K}_0, \mathsf{S}_0$ are inhabited; in
general, for our inductive scheme, inhabitation of closed instances
is computed by the obvious fixed-point marking).  Consider
\emph{element-uniform} interpretations of $R$: those with
$\sem{R}(a,b)(x,y)$ independent of $(x,y)$.  Define $r_\Theta
\subseteq T_0 \times T_0$ as the least relation closed under the
rules read off the instances in $\Theta$: an instance of (trans) at
types $(\sigma, \rho, \chi)$ contributes ``\,$(\sigma,\rho) \in r
\wedge (\rho,\chi) \in r \Rightarrow (\sigma,\chi) \in r$\,''
provided $\sigma, \rho, \chi$ are inhabited (otherwise the instance
is vacuous), and similarly for the other premise schemes; this is a
finite monotone system, so $r_\Theta$ is computable.  We show:
$\mathrm{Inst}_\Theta \folmodels_\Env g(w,w')$ iff $(\lift{w},
\lift{w'}) \in r_\Theta$.  If $(\lift{w}, \lift{w'}) \in r_\Theta$:
by induction on the closure derivation, in \emph{every} model
validating the instances, $\sem{R}$ relates the canonical points of
the corresponding types (inhabitation supplies the element witnesses
needed to apply the guarded instances, and the derivation supplies
the antecedents), so the goal is valid.  If $(\lift{w}, \lift{w'})
\notin r_\Theta$: take the canonical tagged model with the
element-uniform interpretation $\sem{R}(a,b) = \holds$ iff $(\sigma,
\rho) \in r_\Theta$ for some $\sigma, \rho \in T_0$ with
$\sem{\sigma} = a$, $\sem{\rho} = b$ --- in the canonical tagged
model distinct closed types in $T_0$ have distinct denotations
(distinguishable by $\mathrm{tr}$ and cardinality as above), so this
is well-defined --- and $\emptyset$ otherwise.  Every instance in
$\Theta$ is then valid: its types lie in $T_0$, and $r_\Theta$ is
closed under exactly the rule the instance expresses (vacuity
matching the emptiness side conditions).  The goal fails since
$(\lift{w}, \lift{w'}) \notin r_\Theta$.

\emph{Conclusion.}  Suppose an algorithm as in the statement
existed.  Given ground CL terms $w, w'$, compute $\Theta \coloneqq
M(\Env_{\mathrm{CL}}, P_{\mathrm{CL}}, g(w,w'))$ and decide
$\mathrm{Inst}_\Theta \folmodels g(w,w')$ by Claim~2.  If it holds,
then (instances being consequences, \cref{lem:inst-valid})
$P_{\mathrm{CL}} \folmodels g(w,w')$, so $w \to^{*} w'$ by Claim~1.
If it fails, then by the assumed property of $M$ (contrapositive),
$P_{\mathrm{CL}} \not\folmodels g(w,w')$, so $w \not\to^{*} w'$ by
Claim~1.  This decides $w \to^{*} w'$ --- contradiction.
\end{proof}

\begin{remark}
The theorem justifies the two design decisions of this section:
instance selection must be heuristic and bounded
(\cref{def:saturation}), and completeness should be recovered, when
desired, by the \emph{keep} policy
(\cref{prop:mono-conservative}) rather than by cleverer instance
selection.  This mirrors the design of Bobot--Paskevich's
\textsf{Dis} transformation, which keeps the polymorphic original
alongside the instances and is complete for exactly that reason
\cite{BobotPaskevich2011}, and Sledgehammer's heuristic
monomorphisation, which drops generics and accepts incompleteness
\cite[\S5.6]{Blanchette2016}.  The guard-erasure results of
\cref{sec:erasure} apply to the \emph{drop} policy: erasability is a
property of fully monomorphic problems.
\end{remark}
